\documentclass[../DoAn.tex]{subfiles}
\begin{document}

Chương 1 đã giới thiệu bối cảnh, mục tiêu và định hướng giải pháp của đề tài.
Chương này tiến hành khảo sát hiện trạng các công cụ quan sát mô phỏng giao
thông hiện có, phân tích yêu cầu của hệ thống và mô tả các chức năng dưới dạng
biểu đồ use case. Nội dung gồm bốn phần: khảo sát hiện trạng (mục \ref{section:2.1}),
tổng quan chức năng qua biểu đồ use case (mục \ref{section:2.2}), đặc tả chi tiết
các use case quan trọng (mục \ref{section:2.3}) và các yêu cầu phi chức năng
(mục \ref{section:2.4}).

\section{Khảo sát hiện trạng}
\label{section:2.1}

Việc khảo sát hiện trạng được thực hiện dựa trên ba nguồn: nhu cầu của người dùng
khi quan sát mô phỏng giao thông, các công cụ quan sát đi kèm trình mô phỏng, và
các ứng dụng trực quan hóa giao thông ba chiều tương tự. Mục tiêu là xác định các
hạn chế còn tồn tại, từ đó rút ra những tính năng cần phát triển.

\subsection{Công cụ quan sát đi kèm trình mô phỏng}

\gls{SUMO} cung cấp sẵn công cụ quan sát \texttt{sumo-gui} với giao diện hai
chiều theo góc nhìn từ trên cao. Công cụ này hiển thị đầy đủ mạng đường, phương
tiện, làn và trạng thái đèn tín hiệu, đồng thời cho phép phóng to, thu nhỏ và
truy vấn thuộc tính của từng đối tượng. Ưu điểm của \texttt{sumo-gui} là bám sát
dữ liệu mô phỏng, nhẹ và ổn định. Hình \ref{fig:sumo-gui} minh họa giao diện hai
chiều của công cụ này. Tuy nhiên, do bản chất là góc nhìn từ trên cao,
công cụ này phù hợp cho việc giám sát tổng thể nhưng khó tái hiện quan hệ không
gian theo chiều cao và không cung cấp góc nhìn của người trực tiếp tham gia giao
thông.

\begin{figure}[H]
	\centering
	\includegraphics[width=0.8\linewidth]{Hinhve/hinh2.1-sumo-gui.png}
	\caption{Giao diện hai chiều nhìn từ trên cao của \texttt{sumo-gui}}
	\label{fig:sumo-gui}
\end{figure}

\subsection{Các ứng dụng trực quan hóa giao thông ba chiều}

Trong nhóm các giải pháp hiển thị ba chiều, có thể kể đến hai hướng tiêu biểu.
Hướng thứ nhất là các nền tảng mô phỏng phục vụ xe tự hành như CARLA, cung cấp
môi trường ba chiều chân thực cùng mô hình cảm biến. Tuy nhiên, các nền tảng này
yêu cầu tài nguyên lớn, tập trung vào mô phỏng cảm biến cho một số ít phương tiện
thay vì hiển thị dòng giao thông quy mô lớn từ một trình mô phỏng vi mô. Hướng
thứ hai là các công cụ chuyển dữ liệu mô phỏng sang môi trường ba chiều để phát
lại; các công cụ này giúp quan sát trực quan hơn nhưng thường phụ thuộc vào bản
đồ dựng sẵn, ít hỗ trợ thay đổi kịch bản nhanh và hiếm khi cho phép người dùng
can thiệp, điều khiển trực tiếp một phương tiện trong cảnh.

\subsection{So sánh và rút ra yêu cầu}

Bảng \ref{tab:khaosat-sosanh} so sánh các công cụ hiện có với giải pháp đề xuất
của đồ án theo các tiêu chí quan trọng đối với nhu cầu quan sát và phân tích.

\begin{table}[H]
	\centering
	\caption{So sánh các công cụ quan sát mô phỏng giao thông}
	\label{tab:khaosat-sosanh}
	\begin{tabular}{p{0.30\textwidth} c c c}
		\hline
		\textbf{Tiêu chí} & \textbf{sumo-gui} & \textbf{CARLA} & \textbf{Đồ án} \\
		\hline
		Hiển thị ba chiều & Không & Có & Có \\
		Dữ liệu từ mô phỏng vi mô & Có & Hạn chế & Có \\
		Dựng cảnh tự động theo kịch bản & Có & Hạn chế & Có \\
		Hiển thị quy mô lớn nhiều đối tượng & Có & Hạn chế & Có \\
		Góc nhìn người tham gia giao thông & Không & Có & Có \\
		Tương tác lái xe trong cảnh & Không & Có & Có \\
		Yêu cầu tài nguyên & Thấp & Cao & Trung bình \\
		\hline
	\end{tabular}
\end{table}

Từ bảng so sánh, có thể thấy khoảng trống chính nằm ở khả năng kết hợp đồng thời
ba yếu tố: hiển thị ba chiều một lượng lớn phương tiện mà vẫn giữ được mức chi
tiết tới từng xe của mô phỏng vi mô, dựng cảnh tự động theo kịch bản, và cho phép
tương tác theo góc nhìn người tham gia giao thông với yêu cầu tài nguyên vừa
phải. Trên cơ sở đó, hệ thống cần phát triển các tính năng
quan trọng gồm: dựng cảnh ba chiều tự động từ dữ liệu mạng đường; hiển thị phương
tiện, đèn tín hiệu và người đi bộ theo thời gian; cung cấp các thao tác quan sát
như điều khiển camera, tạm dừng và điều chỉnh tốc độ; và cho phép người dùng
chiếm quyền điều khiển, lái thử một phương tiện trong cảnh.

\section{Tổng quan chức năng}
\label{section:2.2}

\subsection{Biểu đồ use case tổng quát}
\label{subsection:2.2.1}

Hệ thống có một tác nhân chính là \textbf{Người dùng}. Đây là người chuẩn bị kịch
bản, khởi chạy và quan sát mô phỏng, thực hiện các thao tác điều khiển quá trình
chạy và xuất kết quả; ngoài ra, cùng người dùng này có thể chiếm quyền điều khiển
một phương tiện để lái thủ công. Hệ thống vận hành theo mô hình một người dùng
cho mỗi phiên mô phỏng, do đó vai trò ``lái xe'' không phải một người riêng biệt
mà chỉ là một chế độ tương tác của chính người dùng đó. Bên cạnh đó, \textbf{trình
mô phỏng \gls{SUMO}} đóng vai trò một tác nhân phụ (hệ thống ngoài) chịu trách
nhiệm sinh dữ liệu mô phỏng vi mô và thực thi các lệnh điều khiển nhận được.

Các use case của hệ thống được nhóm thành năm nhóm chức năng: chuẩn bị kịch bản
và bản đồ, vận hành mô phỏng, quan sát và điều khiển, lái xe tương tác, và kết
quả. Lưu ý việc hiển thị phương tiện, đèn tín hiệu và người đi bộ là hành vi tự
động của hệ thống xảy ra trong quá trình chạy mô phỏng (nhóm B), không phải thao
tác độc lập của người dùng; chúng được mô hình hóa là quan hệ \textit{«include»}
bên trong UC6 và UC7. Biểu đồ use case tổng quát ở Hình \ref{fig:uc-tongquat}
thể hiện quan hệ giữa tác nhân với năm nhóm chức năng này.

\begin{figure}[H]
	\centering
	\begin{tikzpicture}
		% Tac nhan
		\pic at (0,-2) {actor}; \coordinate (op) at (0,-2);
		\node[align=center, font=\small] at (0,-3.1) {Người dùng};
		\pic at (11,-1.6) {actor}; \coordinate (sumo) at (11,-1.6);
		\node[align=center, font=\small] at (11,-2.7) {SUMO\\(hệ thống ngoài)};
		% Bien he thong
		\node[draw, rounded corners, minimum width=6.5cm, minimum height=5.6cm,
		      label=above:{\small \textbf{Hệ thống mô phỏng--hiển thị 3D}}] (sys) at (5.5,-2) {};
		% Nhom chuc nang
		\node[ucbox] (g1) at (5.5,0.2) {A. Chuẩn bị kịch bản \& bản đồ};
		\node[ucbox] (g2) at (5.5,-1.1) {B. Vận hành mô phỏng};
		\node[ucbox] (g3) at (5.5,-2.4) {C. Quan sát \& điều khiển};
		\node[ucbox] (g4) at (5.5,-3.7) {D. Lái xe tương tác};
		\node[ucbox] (g5) at (5.5,-4.9) {E. Kết quả \& lưu trữ};
		% Lien ket
		\draw (op) -- (g1.west);
		\draw (op) -- (g2.west);
		\draw (op) -- (g3.west);
		\draw (op) -- (g4.west);
		\draw (op) -- (g5.west);
		\draw (g2.east) -- (sumo);
		\draw (g4.east) -- (sumo);
	\end{tikzpicture}
	\caption{Biểu đồ use case tổng quát (bản nháp)}
	\label{fig:uc-tongquat}
\end{figure}

Bảng \ref{tab:uc-list} liệt kê toàn bộ các use case của hệ thống theo năm nhóm,
làm cơ sở cho các biểu đồ phân rã ở các mục tiếp theo.

\begin{table}[H]
	\centering
	\caption{Danh sách use case của hệ thống}
	\label{tab:uc-list}
	\begin{tabular}{c p{0.55\textwidth} c}
		\hline
		\textbf{Mã} & \textbf{Use case} & \textbf{Tác nhân} \\
		\hline
		\multicolumn{3}{l}{\textit{A. Chuẩn bị kịch bản \& bản đồ}} \\
		UC1 & Tạo bản đồ từ lưới mê cung & Người dùng \\
		UC2 & Nhập bản đồ từ OpenStreetMap & Người dùng \\
		UC3 & Sinh tuyến đường cho phương tiện & Người dùng \\
		UC4 & Sinh tuyến đường cho người đi bộ & Người dùng \\
		UC5 & Cấu hình tham số mô phỏng & Người dùng \\
		\multicolumn{3}{l}{\textit{B. Vận hành mô phỏng}} \\
		UC6 & Chạy mô phỏng thời gian thực & Người dùng \\
		UC7 & Chạy mô phỏng chế độ phát lại & Người dùng \\
		UC8 & Chạy kèm giao diện sumo-gui & Người dùng \\
		UC9 & Chạy kịch bản định tuyến giao hàng & Người dùng \\
		\multicolumn{3}{l}{\textit{C. Quan sát \& điều khiển}} \\
		UC10 & Điều khiển camera quan sát & Người dùng \\
		UC11 & Tạm dừng / tiếp tục mô phỏng & Người dùng \\
		UC12 & Điều chỉnh tốc độ mô phỏng & Người dùng \\
		UC13 & Lọc đối tượng hiển thị theo vùng & Người dùng \\
		\multicolumn{3}{l}{\textit{D. Lái xe tương tác}} \\
		UC14 & Chiếm quyền điều khiển và lái xe & Người dùng \\
		UC15 & Trả quyền điều khiển phương tiện & Người dùng \\
		\multicolumn{3}{l}{\textit{E. Kết quả \& lưu trữ}} \\
		UC16 & Xử lý va chạm sinh xác xe & Người dùng \\
		UC17 & Ghi và lưu trữ phiên mô phỏng & Người dùng \\
		\hline
	\end{tabular}
\end{table}

\subsection{Biểu đồ use case phân rã nhóm Chuẩn bị kịch bản và bản đồ}
\label{subsection:2.2.2}

Nhóm này gồm các use case phục vụ tạo dữ liệu đầu vào cho mô phỏng. Người vận
hành tạo mạng đường (từ lưới mê cung hoặc từ bản đồ OpenStreetMap), sinh tuyến
đường cho phương tiện và người đi bộ, và cấu hình các tham số mô phỏng. Hình
\ref{fig:uc-nhomA} thể hiện các use case của nhóm này.

\begin{figure}[H]
	\centering
	\begin{tikzpicture}
		\pic at (0,-2) {actor}; \coordinate (op) at (0,-2);
		\node[align=center, font=\small] at (0,-3.1) {Người dùng};
		\node[uc] (u1) at (5,0)    {UC1. Tạo bản đồ\\từ mê cung};
		\node[uc] (u2) at (5,-1.5) {UC2. Nhập bản đồ\\OpenStreetMap};
		\node[uc] (u3) at (5,-3)   {UC3. Sinh tuyến\\cho phương tiện};
		\node[uc] (u4) at (5,-4.5) {UC4. Sinh tuyến\\cho người đi bộ};
		\node[uc] (u5) at (5,-6)   {UC5. Cấu hình\\tham số mô phỏng};
		\foreach \u in {u1,u2,u3,u4,u5} \draw (op) -- (\u.west);
	\end{tikzpicture}
	\caption{Use case nhóm Chuẩn bị kịch bản và bản đồ (bản nháp)}
	\label{fig:uc-nhomA}
\end{figure}

\subsection{Biểu đồ use case phân rã nhóm Vận hành mô phỏng}
\label{subsection:2.2.3}

Nhóm này gồm các use case khởi chạy mô phỏng theo các chế độ khác nhau: thời gian
thực, phát lại, kèm giao diện \texttt{sumo-gui} để đối chiếu, và kịch bản định
tuyến giao hàng. Các use case này tương tác với trình mô phỏng \gls{SUMO}. UC6 và
UC7 tự động bao gồm việc hiển thị phương tiện, đèn tín hiệu và người đi bộ dưới
dạng quan hệ \textit{«include»}. Hình \ref{fig:uc-nhomB} thể hiện nhóm chức năng
này.

\begin{figure}[H]
	\centering
	\begin{tikzpicture}
		\pic at (0,-1.6) {actor}; \coordinate (op) at (0,-1.6);
		\node[align=center, font=\small] at (0,-2.7) {Người dùng};
		\pic at (11,-1.6) {actor}; \coordinate (sumo) at (11,-1.6);
		\node[align=center, font=\small] at (11,-2.7) {SUMO};
		\node[uc] (u6) at (4.5,0)    {UC6. Chạy thời\\gian thực};
		\node[uc] (u7) at (4.5,-1.5) {UC7. Chạy chế độ\\phát lại};
		\node[uc] (u8) at (4.5,-3)   {UC8. Chạy kèm\\sumo-gui};
		\node[uc] (u9) at (4.5,-4.5) {UC9. Kịch bản\\giao hàng};
		% include sub-use-cases
		\node[uc, minimum width=2.6cm, font=\tiny] (inc1) at (8.5,0.6)  {Hiển thị phương tiện};
		\node[uc, minimum width=2.6cm, font=\tiny] (inc2) at (8.5,-0.4) {Hiển thị đèn tín hiệu};
		\node[uc, minimum width=2.6cm, font=\tiny] (inc3) at (8.5,-1.4) {Hiển thị người đi bộ};
		\foreach \u in {u6,u7,u8,u9} {\draw (op) -- (\u.west); \draw (\u.east) -- (sumo);}
		\draw[dashed,-{Latex}] (u6.east) -- node[above, font=\tiny] {«include»} (inc1.west);
		\draw[dashed,-{Latex}] (u6.east) -- (inc2.west);
		\draw[dashed,-{Latex}] (u7.east) -- (inc3.west);
	\end{tikzpicture}
	\caption{Use case nhóm Vận hành mô phỏng (bản nháp)}
	\label{fig:uc-nhomB}
\end{figure}

\subsection{Biểu đồ use case phân rã nhóm Quan sát và trực quan hóa}
\label{subsection:2.2.4}

Nhóm này gồm các use case hiển thị đối tượng và quan sát cảnh: hiển thị phương
tiện, đèn tín hiệu, người đi bộ và điều khiển camera. Hình \ref{fig:uc-nhomC} thể
hiện nhóm chức năng này.

\begin{figure}[H]
	\centering
	\begin{tikzpicture}
		\pic at (0,-1.6) {actor}; \coordinate (op) at (0,-1.6);
		\node[align=center, font=\small] at (0,-2.7) {Người dùng};
		\node[uc] (u10) at (5,0)    {UC10. Hiển thị\\phương tiện};
		\node[uc] (u11) at (5,-1.5) {UC11. Hiển thị\\đèn tín hiệu};
		\node[uc] (u12) at (5,-3)   {UC12. Hiển thị\\người đi bộ};
		\node[uc] (u13) at (5,-4.5) {UC13. Điều khiển\\camera};
		\foreach \u in {u10,u11,u12,u13} \draw (op) -- (\u.west);
	\end{tikzpicture}
	\caption{Use case nhóm Quan sát và trực quan hóa 3D (bản nháp)}
	\label{fig:uc-nhomC}
\end{figure}

\subsection{Biểu đồ use case phân rã nhóm Lái xe tương tác}
\label{subsection:2.2.5}

Nhóm này gồm các use case điều khiển quá trình chạy và tương tác trực tiếp với
phương tiện. Người dùng tạm dừng/tiếp tục, điều chỉnh tốc độ và lọc đối tượng
hiển thị; khi chiếm quyền một phương tiện, người dùng chuyển sang chế độ lái để
điều khiển và sau đó trả quyền. Hình \ref{fig:uc-nhomD} thể hiện nhóm này.

\begin{figure}[H]
	\centering
	\begin{tikzpicture}
		\pic at (0,-2.8) {actor}; \coordinate (op) at (0,-2.8);
		\node[align=center, font=\small] at (0,-3.9) {Người dùng};
		\pic at (10,-2.5) {actor}; \coordinate (sumo) at (10,-2.5);
		\node[align=center, font=\small] at (10,-3.6) {SUMO};
		\node[uc] (u14) at (5,0.2)  {UC14. Tạm dừng/\\tiếp tục};
		\node[uc] (u15) at (5,-1.3) {UC15. Điều chỉnh\\tốc độ};
		\node[uc] (u16) at (5,-2.8) {UC16. Lọc đối\\tượng hiển thị};
		\node[uc] (u17) at (5,-4.3) {UC17. Chiếm quyền\\và lái xe};
		\node[uc] (u18) at (5,-5.8) {UC18. Trả quyền\\điều khiển};
		\foreach \u in {u14,u15,u16,u17,u18} \draw (op) -- (\u.west);
		\draw (u14.east) -- (sumo);
		\draw (u17.east) -- (sumo);
		\draw (u18.east) -- (sumo);
	\end{tikzpicture}
	\caption{Use case nhóm Lái xe tương tác (bản nháp)}
	\label{fig:uc-nhomD}
\end{figure}

\subsection{Biểu đồ use case phân rã nhóm Kết quả và kết thúc}
\label{subsection:2.2.6}

Nhóm này gồm use case xử lý va chạm sinh xác xe (kết quả của tương tác lái xe) và
use case lưu trữ, kết thúc phiên mô phỏng. Hình \ref{fig:uc-nhomE} thể hiện nhóm
chức năng này.

\begin{figure}[H]
	\centering
	\begin{tikzpicture}
		\pic at (0,-1.5) {actor}; \coordinate (op) at (0,-1.5);
		\node[align=center, font=\small] at (0,-2.6) {Người dùng};
		\node[uc] (u19) at (5,0)    {UC19. Xử lý va chạm\\sinh xác xe};
		\node[uc] (u20) at (5,-3)   {UC20. Lưu trữ \&\\kết thúc phiên};
		\draw (op) -- (u19.west);
		\draw (op) -- (u20.west);
	\end{tikzpicture}
	\caption{Use case nhóm Kết quả và kết thúc (bản nháp)}
	\label{fig:uc-nhomE}
\end{figure}

\section{Đặc tả chức năng}
\label{section:2.3}

Phần này đặc tả chi tiết sáu use case quan trọng nhất, được lựa chọn vì nằm trên
luồng nghiệp vụ chính và có nhiều luồng sự kiện phát sinh. Các use case còn lại
được đặc tả đầy đủ trong Phụ lục B.

\subsection{Đặc tả use case Chạy mô phỏng thời gian thực}

\begin{table}[H]
	\centering
	\caption{Đặc tả use case UC6 -- Chạy mô phỏng thời gian thực}
	\begin{tabular}{p{0.22\textwidth} p{0.70\textwidth}}
		\hline
		\textbf{Tên} & Chạy mô phỏng thời gian thực \\
		\textbf{Tác nhân} & Người dùng, SUMO \\
		\textbf{Mô tả} & Hệ thống chạy mô phỏng và truyền trạng thái từng bước sang ứng dụng hiển thị để cập nhật tức thời. \\
		\textbf{Tiền điều kiện} & Đã có kịch bản mô phỏng hợp lệ; ứng dụng hiển thị đã kết nối. \\
		\textbf{Luồng chính} & (1) Người dùng chọn chế độ thời gian thực và khởi chạy. (2) Hệ thống khởi tạo \gls{SUMO} và mở các kênh truyền. (3) Mỗi bước, hệ thống lấy trạng thái phương tiện, đèn, người đi bộ. (4) Hệ thống đóng gói và gửi sang ứng dụng hiển thị. (5) Ứng dụng hiển thị cập nhật cảnh. (6) Lặp lại đến khi hết mô phỏng hoặc có lệnh dừng. \\
		\textbf{Luồng phụ} & (3a) Gói tin bị phân mảnh trên đường truyền: phía nhận đệm dữ liệu đến khi đủ một bản tin. (2a) Không kết nối được ứng dụng hiển thị: hệ thống thử lại; nếu thất bại thì dừng. \\
		\textbf{Hậu điều kiện} & Trạng thái mô phỏng được hiển thị theo thời gian; dữ liệu phiên được ghi lại. \\
		\hline
	\end{tabular}
\end{table}

\subsection{Đặc tả use case Sinh tuyến đường cho phương tiện}

\begin{table}[H]
	\centering
	\caption{Đặc tả use case UC3 -- Sinh tuyến đường cho phương tiện}
	\begin{tabular}{p{0.22\textwidth} p{0.70\textwidth}}
		\hline
		\textbf{Tên} & Sinh tuyến đường cho phương tiện \\
		\textbf{Tác nhân} & Người dùng \\
		\textbf{Mô tả} & Hệ thống sinh tập tuyến đường cho phương tiện dựa trên cách chọn cặp điểm đầu--cuối và lọc các cặp không liên thông. \\
		\textbf{Tiền điều kiện} & Đã có mạng đường hợp lệ. \\
		\textbf{Luồng chính} & (1) Người dùng chọn kiểu phân chia cặp điểm đầu--cuối và số lượng. (2) Hệ thống sinh các cặp theo phân bố tọa độ. (3) Hệ thống kiểm tra tính liên thông của từng cặp. (4) Hệ thống ghi các tuyến hợp lệ ra tệp tuyến đường. \\
		\textbf{Luồng phụ} & (3a) Cặp điểm không có đường đi: hệ thống loại bỏ cặp đó và tiếp tục. (4a) Không còn cặp hợp lệ: hệ thống báo lỗi và yêu cầu cấu hình lại. \\
		\textbf{Hậu điều kiện} & Tệp tuyến đường được tạo, sẵn sàng cho mô phỏng. \\
		\hline
	\end{tabular}
\end{table}

\subsection{Đặc tả use case Chiếm quyền điều khiển và lái xe}

\begin{table}[H]
	\centering
	\caption{Đặc tả use case UC17 -- Chiếm quyền điều khiển và lái xe}
	\begin{tabular}{p{0.22\textwidth} p{0.70\textwidth}}
		\hline
		\textbf{Tên} & Chiếm quyền điều khiển và lái xe \\
		\textbf{Tác nhân} & Người dùng, SUMO \\
		\textbf{Mô tả} & Người dùng chọn một phương tiện đang do \gls{SUMO} điều khiển và chuyển nó sang chế độ lái thủ công bằng vật lý. \\
		\textbf{Tiền điều kiện} & Mô phỏng đang chạy; tồn tại phương tiện có thể chiếm quyền. \\
		\textbf{Luồng chính} & (1) Người dùng chọn phương tiện và yêu cầu chiếm quyền. (2) Hệ thống chuyển phương tiện sang chế độ điều khiển bởi người dùng và bật vật lý. (3) Người dùng điều khiển bằng bàn phím. (4) Hệ thống phản ánh vị trí xe sang \gls{SUMO} để các xe khác phản ứng. \\
		\textbf{Luồng phụ} & (3a) Xe đang lái va chạm với xe khác: hệ thống đánh dấu trạng thái va chạm để xử lý khi trả quyền (xem UC19). \\
		\textbf{Hậu điều kiện} & Phương tiện ở trạng thái do người dùng điều khiển. \\
		\hline
	\end{tabular}
\end{table}

Do use case này có luồng trạng thái phức tạp, biểu đồ hoạt động ở Hình
\ref{fig:act-uc17} mô tả luồng chiếm quyền và lái xe.

\begin{figure}[H]
	\centering
	\begin{tikzpicture}[node distance=0.9cm and 0cm]
		\node[circle, fill=black, inner sep=2.5pt] (s) {};
		\node[flow, below=of s] (a1) {Chọn xe \& yêu cầu chiếm quyền};
		\node[flow, below=of a1] (a2) {Chuyển sang chế độ người dùng,\\bật vật lý};
		\node[flow, below=of a2] (a3) {Người dùng điều khiển bằng bàn phím};
		\node[flow, below=of a3] (a4) {Phản ánh vị trí xe sang SUMO};
		\node[decision, below=of a4] (d1) {Có va chạm?};
		\node[flow, below=of d1] (a5) {Đánh dấu trạng thái va chạm};
		\node[circle, draw, inner sep=3pt, below=1.1cm of a5] (e) {\footnotesize kết};
		\draw[-{Latex}] (s) -- (a1);
		\draw[-{Latex}] (a1) -- (a2);
		\draw[-{Latex}] (a2) -- (a3);
		\draw[-{Latex}] (a3) -- (a4);
		\draw[-{Latex}] (a4) -- (d1);
		\draw[-{Latex}] (d1) -- node[right, font=\small] {Có} (a5);
		\draw[-{Latex}] (a5) -- (e);
		\draw[-{Latex}] (d1.west) -- ++(-2.6,0) node[above, font=\small] {Không} |- (e.west);
	\end{tikzpicture}
	\caption{Biểu đồ hoạt động use case Chiếm quyền và lái xe (bản nháp)}
	\label{fig:act-uc17}
\end{figure}

\subsection{Đặc tả use case Trả quyền điều khiển phương tiện}

\begin{table}[H]
	\centering
	\caption{Đặc tả use case UC18 -- Trả quyền điều khiển phương tiện}
	\begin{tabular}{p{0.22\textwidth} p{0.70\textwidth}}
		\hline
		\textbf{Tên} & Trả quyền điều khiển phương tiện \\
		\textbf{Tác nhân} & Người dùng, SUMO \\
		\textbf{Mô tả} & Người dùng thả một phương tiện đang lái; hệ thống đưa xe trở lại làn đường để \gls{SUMO} tiếp tục điều khiển, hoặc loại bỏ xe nếu không thể đưa về làn. \\
		\textbf{Tiền điều kiện} & Phương tiện đang ở trạng thái do người dùng điều khiển. \\
		\textbf{Luồng chính} & (1) Người dùng yêu cầu trả quyền. (2) Hệ thống tìm làn đường gần nhất với vị trí hiện tại của xe. (3) Nếu khoảng cách trong ngưỡng cho phép, hệ thống đưa xe về làn, cấp tuyến mới và trả quyền cho \gls{SUMO}. (4) Phương tiện tiếp tục di chuyển tự động. \\
		\textbf{Luồng phụ} & (2a) Xe ở trạng thái va chạm: chuyển sang xử lý xác xe thay vì trả về \gls{SUMO} (UC19). (3a) Vị trí xe quá xa mọi làn đường: hệ thống loại bỏ phương tiện. \\
		\textbf{Hậu điều kiện} & Phương tiện được trả về \gls{SUMO} hoặc bị loại bỏ. \\
		\hline
	\end{tabular}
\end{table}

Hình \ref{fig:act-uc18} mô tả luồng hoạt động của use case trả quyền điều khiển.

\begin{figure}[H]
	\centering
	\begin{tikzpicture}[node distance=0.9cm and 0cm]
		\node[circle, fill=black, inner sep=2.5pt] (s) {};
		\node[flow, below=of s] (a1) {Người dùng yêu cầu trả quyền};
		\node[decision, below=of a1] (d0) {Xe bị\\va chạm?};
		\node[flow, below=of d0] (a2) {Tìm làn đường gần nhất};
		\node[decision, below=of a2] (d1) {Trong ngưỡng\\khoảng cách?};
		\node[flow, below=of d1] (a3) {Đưa về làn, cấp tuyến,\\trả quyền cho SUMO};
		\node[flow, right=3.2cm of d1] (a4) {Loại bỏ phương tiện};
		\node[flow, right=3.2cm of d0] (a5) {Chuyển xử lý xác xe (UC19)};
		\node[circle, draw, inner sep=3pt, below=1.0cm of a3] (e) {\footnotesize kết};
		\draw[-{Latex}] (s) -- (a1);
		\draw[-{Latex}] (a1) -- (d0);
		\draw[-{Latex}] (d0) -- node[left, font=\small] {Không} (a2);
		\draw[-{Latex}] (d0) -- node[above, font=\small] {Có} (a5);
		\draw[-{Latex}] (a2) -- (d1);
		\draw[-{Latex}] (d1) -- node[left, font=\small] {Có} (a3);
		\draw[-{Latex}] (d1) -- node[above, font=\small] {Không} (a4);
		\draw[-{Latex}] (a3) -- (e);
		\draw[-{Latex}] (a4) |- (e);
		\draw[-{Latex}] (a5) |- (e);
	\end{tikzpicture}
	\caption{Biểu đồ hoạt động use case Trả quyền điều khiển (bản nháp)}
	\label{fig:act-uc18}
\end{figure}

\subsection{Đặc tả use case Tạm dừng và tiếp tục mô phỏng}

\begin{table}[H]
	\centering
	\caption{Đặc tả use case UC14 -- Tạm dừng và tiếp tục mô phỏng}
	\begin{tabular}{p{0.22\textwidth} p{0.70\textwidth}}
		\hline
		\textbf{Tên} & Tạm dừng và tiếp tục mô phỏng \\
		\textbf{Tác nhân} & Người dùng, SUMO \\
		\textbf{Mô tả} & Người dùng tạm dừng hoặc tiếp tục quá trình mô phỏng đang chạy để quan sát kỹ một thời điểm. \\
		\textbf{Tiền điều kiện} & Mô phỏng đang chạy ở chế độ thời gian thực. \\
		\textbf{Luồng chính} & (1) Người dùng gửi lệnh tạm dừng. (2) Hệ thống dừng tiến trình bước mô phỏng nhưng giữ nguyên trạng thái. (3) Người dùng gửi lệnh tiếp tục. (4) Hệ thống tiếp tục bước mô phỏng từ trạng thái đã giữ. \\
		\textbf{Luồng phụ} & (2a) Khi tạm dừng, các đối tượng động (kể cả người đi bộ) phải đứng yên, không tiếp tục chuyển động tại chỗ. \\
		\textbf{Hậu điều kiện} & Mô phỏng ở trạng thái tạm dừng hoặc đang chạy theo lệnh. \\
		\hline
	\end{tabular}
\end{table}

\subsection{Đặc tả use case Nhập bản đồ từ OpenStreetMap}

\begin{table}[H]
	\centering
	\caption{Đặc tả use case UC2 -- Nhập bản đồ từ OpenStreetMap}
	\begin{tabular}{p{0.22\textwidth} p{0.70\textwidth}}
		\hline
		\textbf{Tên} & Nhập bản đồ từ OpenStreetMap \\
		\textbf{Tác nhân} & Người dùng \\
		\textbf{Mô tả} & Hệ thống chuyển một tệp bản đồ OpenStreetMap thành mạng đường và kịch bản mô phỏng. \\
		\textbf{Tiền điều kiện} & Người dùng có tệp bản đồ OpenStreetMap hợp lệ. \\
		\textbf{Luồng chính} & (1) Người dùng chọn tệp bản đồ. (2) Hệ thống chuyển đổi thành mạng đường. (3) Hệ thống sinh tuyến cho phương tiện và người đi bộ. (4) Hệ thống tạo tệp cấu hình mô phỏng. \\
		\textbf{Luồng phụ} & (2a) Bản đồ lỗi hoặc không dựng được mạng: hệ thống báo lỗi và dừng quá trình. \\
		\textbf{Hậu điều kiện} & Kịch bản mô phỏng từ bản đồ thực tế được tạo, sẵn sàng chạy. \\
		\hline
	\end{tabular}
\end{table}

\section{Yêu cầu phi chức năng}
\label{section:2.4}

Bên cạnh các yêu cầu chức năng, hệ thống cần đáp ứng một số yêu cầu phi chức năng
quan trọng.

\textbf{Hiệu năng.} Hệ thống phải duy trì tốc độ khung hình ổn định khi số lượng
phương tiện lớn. Để đạt được điều này, ứng dụng hiển thị áp dụng các kỹ thuật như
tái sử dụng đối tượng, lọc đối tượng ngoài vùng quan sát và giảm số lệnh vẽ. Việc
truyền dữ liệu theo từng bước cũng cần kiểm soát kích thước gói tin để không gây
nghẽn.

\textbf{Độ tin cậy.} Kênh truyền dữ liệu phải chịu được hiện tượng phân mảnh và
gộp gói của giao thức truyền theo luồng, thông qua cơ chế tách khung bản tin. Khi
kết thúc, hệ thống phải đóng kết nối và ghi đầy đủ kết quả, tránh mất dữ liệu hay
rò rỉ tài nguyên.

\textbf{Tính mở rộng.} Dữ liệu trao đổi giữa phần mô phỏng và phần hiển thị cần
được chuẩn hóa theo một định dạng độc lập với nền tảng hiển thị, để có thể thay
thế hoặc bổ sung nền tảng hiển thị khác mà không phải thay đổi phần mô phỏng.

\textbf{Tính dễ dùng.} Người dùng cấu hình và khởi chạy mô phỏng thông qua
giao diện khởi chạy, không cần thao tác trực tiếp với dòng lệnh phức tạp. Các
thao tác quan sát và điều khiển trong khi chạy được cung cấp ngay trên giao diện
ba chiều.

Chương này đã khảo sát hiện trạng các công cụ quan sát mô phỏng giao thông, chỉ
ra khoảng trống về khả năng hiển thị ba chiều có tương tác theo góc nhìn người
tham gia giao thông, và phân tích yêu cầu của hệ thống thông qua hai mươi use case
thuộc năm nhóm chức năng. Các use case quan trọng đã được đặc tả chi tiết kèm
biểu đồ hoạt động cho những luồng phức tạp. Các yêu cầu này là cơ sở để lựa chọn
công nghệ ở Chương 3 và thiết kế hệ thống ở Chương 4.

\end{document}
